\chapter{Baseline Particle Problem}\label{ch:basic}

\section{Physical problem and parameter set}
The baseline separates carrier-flow approximation from particle response.
The carrier is steady and two-dimensional. Particles are spherical, heavy,
and dilute, and their forces do not feed back into the fluid. The principal
channel consists of two connected semicircular bends with nominal inner
and outer radii $R_{\mathrm{in}}=0.2$ and $R_{\mathrm{out}}=0.5$. Its
width is $W=0.3$ and its centerline length is approximately $2.1991$, all
scaled by $L_0$. The Reynolds number is $\Rey=1$.

The baseline particle parameters are $R=1000$ and $d_p=0.025$. The resulting
Stokes number is $\St\simeq0.03472$. The channel wall rule stops the
particle when its center approaches within one radius of a wall. This is
a specified irreversible sticking model, not a prediction of adhesion
chemistry or a resolved impact process. In the full trajectory set, the
largest recorded slip Reynolds number is approximately $0.01221$, which
supports the low-slip-Reynolds-number assumption for these calculations.
The ratio $d_p/W\simeq0.0833$ still warrants caution for near-wall motion.

\section{Flow and force representations}
\subsection{FlowNet}
The flow network maps five input slots to velocity and pressure:
\begin{equation}
 (x,y,R_{\mathrm{in}},R_{\mathrm{out}},\Rey)
 \longmapsto (u,v,p_v).
 \label{eq:flownet}
\end{equation}
Velocity--pressure PINN formulations for incompressible flow provide
related methodological background \citep{jin2021}. The present network
has five hidden layers with 70 tanh units per layer. % Citation-plan C19
 The spatial inputs
are centered and scaled, and the pressure output is scaled with a
channel-based pressure magnitude. The geometry and Reynolds inputs are
fixed within an individual trained model. Their presence does not establish
generalization across untrained geometries or Reynolds numbers.

The objective combines momentum and continuity residuals, wall and end
velocity conditions, a pressure gauge, and a section-flux condition. If
$r_c=\nabla\cdot\vect{u}$ and $\vect{r}_m$ denotes the residual of
\cref{eq:baseline_ns}, the scaled physics term is
\begin{equation}
 \mathcal{L}_{\mathrm{phys}}
 =\left\langle (hr_c)^2\right\rangle
 +\left\langle |h^2\vect{r}_m|^2\right\rangle.
 \label{eq:baseline_loss}
\end{equation}
Here the averages follow the componentwise mean-square convention of the
implementation. The complete objective uses weights 20 for the velocity
boundary term, 1 for the scaled gauge term, and 5 for the normalized flux
term. Finite-element interior values are not used in this objective.

\subsection{ForceNet}
The second network has five hidden layers with 90 tanh units and maps
\begin{equation}
 (u,v,p_v,d_p,x_p,y_p,u_p,v_p)\longmapsto(F_x,F_y).
 \label{eq:forcenet}
\end{equation}
The force scale is the displaced-fluid mass times $U_0^2/L_0$. Under this
convention, the dimensionless Stokes force and particle acceleration are
\begin{equation}
 \vect{F}^{*}=\frac{18}{\Rey d_p^2}
 (\vect{u}-\vect{v}),\qquad
 \frac{d\vect{v}}{dt}=\frac{\vect{F}^{*}}{R}.
 \label{eq:force_scaling}
\end{equation}
The neural approximation is trained on the normalized slip response.
Subtracting its evaluation at zero slip from its evaluation at the actual
slip enforces exactly zero force when $\vect{u}=\vect{v}$. The network
therefore approximates a specified drag relation. Its agreement with that
relation is an approximation test, not discovery of a new constitutive law.

\Cref{fig:basic_architecture} gives a visual summary of both networks.
The role of ForceNet is intentionally kept separate from the carrier-field
approximation so that the two errors can be evaluated independently.

\ResultFigure[.96]{02_baseline_architecture}{Baseline neural architecture.
Each row represents a fully connected MLP; the few displayed nodes are
illustrative, while the numbers above hidden layers give their actual
widths. FlowNet supplies velocity and pressure. ForceNet approximates the
normalized Stokes response and uses a zero-slip subtraction. The first
network is constrained by PDE and boundary residuals; the second by an
algebraic drag residual.}{fig:basic_architecture}

\section{Independent finite-element reference}
The reference solves the same incompressible equations using quadratic
velocity and linear pressure Taylor--Hood elements
\citep{taylor1973,gustafsson2020}. A Picard iteration treats the convective
term, and sparse direct solution is used for the linear subproblems.
The finest channel mesh contains 3,026 triangles, 12,530 velocity degrees
of freedom, and 1,620 pressure degrees of freedom.

Three meshes assess sensitivity. \Cref{tab:mesh} shows that the difference
between the two finest velocity fields is \MeshDifference\%. The finest
reference has a nonlinear residual of approximately
$1.60\times10^{-10}$. Its reported divergence norm is not identically zero;
the mixed discretization enforces incompressibility weakly. The mesh
differences are sensitivity indicators, not a formal uncertainty estimate
or a grid-convergence index \citep{roache1998,celik2008}. % Citation-plan C20

\begin{table}[htbp]\centering\small
\caption{Reference-mesh assessment. Velocity differences are measured
relative to the finest field on the comparison samples.}
\label{tab:mesh}
\begin{tabularx}{\textwidth}{@{}rrrrY@{}}\toprule
Triangles & Velocity DOFs & Pressure DOFs & Difference (\%) & Nonlinear residual\\\midrule
\EvidenceRows{mesh}
\bottomrule\end{tabularx}
\end{table}

The first successful training seed is used for the principal comparison;
it is not selected by its error against the reference. The reported final
evaluation restores completed checkpoints. Consequently, the training
budgets associated with those checkpoints, rather than the defaults of a
later evaluation command, identify the computations that produced the
weights. The computational record in \cref{app:reproducibility} preserves
this distinction.

\section{Analytical plane-flow verification}
Before interpreting curved-channel results, I consider plane Poiseuille
flow in a straight channel with half-width $h$. In the baseline units, the
analytical solution is
\begin{equation}
 u=1-(y/h)^2,\qquad v=0,\qquad
 p_v=-\frac{2x}{h^2}+C.
 \label{eq:poiseuille}
\end{equation}
The finite-element relative velocity error is approximately
$8.19\times10^{-10}$. The PINN velocity error is
\PoiseuilleVelocity\%, and its gauge-aligned pressure error is
\PoiseuillePressure\%. The velocity metric includes both components;
\cref{fig:poiseuille} displays the streamwise profile for readability.
Pressure is evaluated on points distributed along the channel, with each
array's mean removed.

The near-exact finite-element result for this polynomial plane flow is
consistent with the approximation spaces and simple geometry. It should
not be interpreted as the error of the curved-channel reference. The
curved problem requires its own mesh-sensitivity assessment.

\ResultFigure{03_poiseuille}{Analytical verification of the baseline flow
solvers. Panel (a) compares the streamwise velocity profile with the
parabolic solution. Panel (b) compares pressure after constant-gauge
alignment on spatially distributed samples. Lines and markers retain the
stored predictions; no fitted correction is applied.}{fig:poiseuille}

\section{Curved-channel carrier-flow results}
\Cref{fig:curved_fields} compares the principal neural field with the
finite-element field. Both solutions recover the curved high-speed core
and the pressure decrease through the channel. Common color limits are
used for corresponding speed and pressure panels. The discrepancy panels
show that agreement is spatially nonuniform; a small global norm should
not be interpreted as uniformly small local error.

The velocity error is \FlowVelocity\%, and the centered pressure error is
\FlowPressure\% on the independent field-comparison samples. The maximum
absolute velocity discrepancy is approximately $0.05401$. These values
support agreement for the specified boundary-value problem. They do not
establish zero equation residual: the separately recorded raw continuity
and momentum RMS values are approximately $0.1781$ and $1.8296$ in the
baseline equation convention. Those residuals, the field errors, and the
flux diagnostic measure different aspects of the approximation.

\ResultFigure{04_curved_fields}{Curved-channel carrier field at $\Rey=1$.
The speed panels use $|\vect{u}|/U_0$ and identical scales; the pressure
panels use centered viscous-scale pressure and identical scales.
Velocity discrepancy is the Euclidean component error, whereas pressure
discrepancy is signed. The fields are evaluated on the same stored mapped
grid; no spanwise solution is implied.}{fig:curved_fields}

\begin{table}[htbp]\centering\small
\caption{Supported baseline numerical comparisons. The relative errors
refer to their stated sample sets and normalizations.}
\label{tab:basic_errors}
\begin{tabularx}{\textwidth}{@{}p{.31\textwidth}Yr@{}}\toprule
Problem & Quantity & Error\\\midrule
\EvidenceRows{basic_errors}
\bottomrule\end{tabularx}
\end{table}

\section{Held-out drag and field comparisons}
The normalized ForceNet error is \DragError\% on the independent
2,048-sample evaluation used for the reported metric.
\Cref{fig:heldout} presents velocity and pressure parity together with a
separate 1,000-sample drag diagnostic. The plotted drag set and the aggregate
2,048-sample test are distinct; they need not have identical numerical
errors. This distinction avoids assigning one test's error to another
test's points.

The drag comparison supports the network's approximation of the retained
creeping-flow law. It cannot establish the accuracy of neglected lift,
history, lubrication, or wall corrections
\citep{maxey1983,goldman1967,brenner1961}. % Citation-plan C21
 Direct analytical Stokes drag
also remains an essential baseline for judging the value of a learned
approximation of this known relation.

\ResultFigure{05_heldout_agreement}{Independent numerical comparisons.
Velocity components and gauge-aligned pressure are compared with the
finite-element reference. The third panel compares normalized learned
drag with the specified Stokes-slip response on the stored 1,000-sample
diagnostic. The dashed lines indicate equality. The aggregate drag error
in \cref{tab:basic_errors} comes from a separate 2,048-sample test.}
{fig:heldout}

\section{Seed variation, reference sensitivity, and flux}
Two independently trained seeds give velocity errors of 1.510\% and
1.429\%, and pressure errors of 0.863\% and 0.886\%, respectively. Their
mean velocity error is \SeedVelocityMean\%, with a sample standard deviation
of \SeedVelocitySD\ percentage points. Two seeds are sufficient to reveal
some training dependence but are insufficient for a strong statistical
uncertainty statement. No confidence interval is inferred.

\begin{table}[htbp]\centering
\caption{Field errors for the two completed training seeds.}
\label{tab:seeds}
\begin{tabular}{@{}rrr@{}}\toprule
Seed & Velocity error (\%) & Pressure error (\%)\\\midrule
\EvidenceRows{seeds}
\bottomrule\end{tabular}
\end{table}

\ResultFigure{06_reference_and_seeds}{Reference and training sensitivity.
The left panel compares the two coarser finite-element meshes with the
finest reference; the right panel shows the two completed training seeds.
The finest mesh is a comparison reference, not an exact solution. Seed
labels identify the individual observations rather than confidence bounds.}
{fig:reference_seeds}

The maximum section-flux deviation is \FluxDeviation\%.
\Cref{fig:flux} shows the full stored section sequence normalized by the
prescribed flux. The variation is finite and should remain visible even
though the study's 5\% flux criterion is met. Together with boundary RMS
of approximately $0.00420$, this result provides an integrated check on
the velocity field that complements the sample-relative error.

\ResultFigure[.95]{07_flux}{Section flux in the curved baseline, normalized
by the prescribed value $Q_0=4h/3=0.2$. The dashed line is the exact
constant-flux target, and every stored PINN section is shown. The maximum
deviation is \FluxDeviation\%; the curve has not been rescaled to enforce
agreement.}{fig:flux}

\section{Particle-transport interpretation}
The baseline evaluation completes all 119 particle or tracer integrations:
49 use PINN with ForceNet, 49 use
FEM with exact Stokes drag, seven use PINN with exact Stokes drag, and
14 are tracer integrations split equally between the PINN and FEM fields.
These counts include the parameter-study cases; the primary endpoint
comparison concerns seven initial conditions. The reference velocity
interpolator differs from direct finite-element evaluation by
\InterpolationDifference\% on its assessment samples. These are useful
computational checks, but neither alone proves accurate particle outcomes.

Among the seven baseline particle initial positions, six endpoint classes
agree with the reference. The largest width-normalized RMS position
discrepancy is approximately 34.97\%, and the lowest common-time coverage
is approximately 73.55\%. Three initial conditions meet all endpoint,
position-error and common-time criteria; six meet the endpoint criterion
alone. Detailed values are retained in
\cref{tab:path_limits} so that the scope of the carrier-flow findings is
unambiguous. The quantitative conclusions of this chapter concern the
carrier field, drag approximation, and the seven tabulated path comparisons.

This outcome is physically and numerically informative. Bulk field
agreement and a good approximation of a specified drag law are necessary
components, but their combination is not enough to certify wall-sensitive
transport. The justified baseline conclusion is that the carrier and drag
components agree well with their stated references, while the full
trajectory prediction remains an improvement target.

The present comparisons do not isolate a single cause for the trajectory
differences. Carrier-field approximation, ForceNet error, reference
interpolation, ODE integration, and event handling are separate candidates.
The recorded integrator uses relative and absolute tolerances
$2\times10^{-6}$ and $2\times10^{-8}$, with RK45 or Radau chosen according
to the response scale. A factorial comparison holding the integrator and
wall-event rule fixed is needed to attribute the discrepancies.

\section{Chapter conclusions}
The baseline establishes a reproducible carrier-flow and drag foundation
for the subsequent coupling study. Analytical plane flow, independent
curved-channel fields, two completed training seeds, and section flux
provide complementary evidence. The results also set a boundary on the
claim: the present field accuracy does not establish reliable deposition
or all baseline particle endpoints. That boundary becomes more important
when reaction forcing and geometric narrowing are introduced.
